\chapter{Release 2 --- Mise en place de l'infrastructure}

\section*{Introduction}
\addcontentsline{toc}{section}{Introduction}
Cette release couvre la containerisation des microservices et leur déploiement, d'abord sur un cluster Kubernetes local (Minikube), puis sur l'infrastructure cloud AWS. Elle est répartie sur deux sprints.

%======================================================================
% SPRINT 4
%======================================================================
\section{Sprint 4 --- Containerisation et déploiement local}

\subsection{Objectif du sprint}
Containeriser l'ensemble des microservices avec Docker et valider leur déploiement sur un cluster Kubernetes local avec Minikube.

\subsection{Sprint Backlog}

\begin{table}[H]
\centering
\renewcommand{\arraystretch}{1.2}
\begin{tabular}{|c|p{9cm}|c|c|}
\hline
\textbf{ID} & \textbf{Tâche} & \textbf{Priorité} & \textbf{Points} \\
\hline
4.1 & Rédaction des Dockerfiles multi-étapes pour chaque microservice & Must & 5 \\
\hline
4.2 & Rédaction du fichier docker-compose.yml pour l'environnement local & Must & 5 \\
\hline
4.3 & Rédaction des manifests Kubernetes (Deployment, Service, ConfigMap) & Must & 8 \\
\hline
4.4 & Déploiement et validation sur Minikube & Must & 5 \\
\hline
\end{tabular}
\caption{Sprint Backlog --- Sprint 4}
\label{tab:sprint4_backlog}
\end{table}

\subsection{Containerisation avec Docker}

Chaque microservice dispose d'un Dockerfile multi-étapes qui sépare la phase de compilation (build) de la phase d'exécution (runtime), permettant de réduire la taille de l'image finale.

\begin{lstlisting}[language=bash, caption={Exemple de Dockerfile multi-étapes}, label={lst:dockerfile}]
# Etape 1 : Compilation
FROM maven:3.9-eclipse-temurin-21 AS build
WORKDIR /app
COPY pom.xml .
COPY src ./src
RUN mvn clean package -DskipTests

# Etape 2 : Execution
FROM eclipse-temurin:21-jre-alpine
COPY --from=build /app/target/*.jar app.jar
EXPOSE 8081
ENTRYPOINT ["java", "-jar", "app.jar"]
\end{lstlisting}

\subsection{Orchestration locale avec Docker Compose}
Le fichier \texttt{docker-compose.yml} orchestre l'ensemble des services en environnement de développement local : base de données MySQL, les microservices Spring Boot, l'API Gateway, le frontend Angular et le serveur Eureka.

\subsection{Déploiement sur Minikube}
Minikube fournit un cluster Kubernetes local pour valider les manifests avant le déploiement en production. Les manifests sont organisés en quatre répertoires :
\begin{itemize}
    \item \texttt{infra/} : MySQL (StatefulSet avec PersistentVolumeClaim).
    \item \texttt{config/} : ConfigMap et Secrets pour les variables d'environnement.
    \item \texttt{services/} : Deployments et Services pour chaque microservice.
    \item \texttt{frontend/} : Deployment et Service pour l'application Angular.
\end{itemize}

\begin{figure}[H]
    \centering
    \fbox{\parbox{0.85\textwidth}{\centering\vspace{2.5cm}\textbf{[Capture d'écran --- Pods sur Minikube]}\\\texttt{kubectl get pods -n billing-dev}\vspace{2.5cm}}}
    \caption{Pods déployés sur le cluster Minikube}
    \label{fig:minikube_pods}
\end{figure}

%======================================================================
% SPRINT 5
%======================================================================
\section{Sprint 5 --- Déploiement sur le cloud AWS}

\subsection{Objectif du sprint}
Provisionner l'infrastructure cloud AWS et déployer la plateforme en environnement de production.

\subsection{Sprint Backlog}

\begin{table}[H]
\centering
\renewcommand{\arraystretch}{1.2}
\begin{tabular}{|c|p{9cm}|c|c|}
\hline
\textbf{ID} & \textbf{Tâche} & \textbf{Priorité} & \textbf{Points} \\
\hline
5.1 & Provisionnement du cluster EKS via la console et CLI AWS & Must & 8 \\
\hline
5.2 & Création de l'instance RDS MySQL & Must & 5 \\
\hline
5.3 & Configuration du registre ECR pour les images Docker & Must & 3 \\
\hline
5.4 & Adaptation des manifests K8s pour l'environnement cloud & Must & 5 \\
\hline
5.5 & Premier déploiement sur EKS et validation & Must & 5 \\
\hline
\end{tabular}
\caption{Sprint Backlog --- Sprint 5}
\label{tab:sprint5_backlog}
\end{table}

\subsection{Architecture de l'infrastructure AWS}
L'infrastructure de production repose sur les services managés suivants :
\begin{itemize}
    \item \textbf{Amazon EKS :} Cluster Kubernetes managé pour l'orchestration des conteneurs.
    \item \textbf{Amazon RDS :} Instance MySQL managée pour la persistance des données.
    \item \textbf{Amazon ECR :} Registre privé pour le stockage des images Docker.
    \item \textbf{Amazon EC2 :} Instance hébergeant le serveur Jenkins pour le pipeline CI/CD.
\end{itemize}

\begin{figure}[H]
    \centering
    \fbox{\parbox{0.85\textwidth}{\centering\vspace{3cm}\textbf{[Diagramme d'architecture AWS]}\\\textit{EC2 (Jenkins) $\rightarrow$ ECR $\rightarrow$ EKS $\rightarrow$ RDS}\vspace{3cm}}}
    \caption{Architecture de l'infrastructure AWS}
    \label{fig:aws_architecture}
\end{figure}

\subsection{Provisionnement du cluster EKS}
Le cluster EKS a été provisionné via la console AWS et l'outil \texttt{eksctl}. La configuration inclut la création de groupes de n\oe{}uds (node groups) avec les ressources nécessaires pour l'exécution des microservices.

\subsection{Configuration d'Amazon RDS}
Une instance MySQL 8.0 a été créée sur Amazon RDS. Les bases de données suivantes sont provisionnées :
\texttt{db\_authentication}, \texttt{db\_customer}, \texttt{db\_catalog}, \texttt{db\_subscription}, \texttt{db\_boutique}, \texttt{db\_billing}.

\begin{figure}[H]
    \centering
    \fbox{\parbox{0.85\textwidth}{\centering\vspace{2.5cm}\textbf{[Capture d'écran --- Console AWS EKS]}\vspace{2.5cm}}}
    \caption{Cluster EKS sur la console AWS}
    \label{fig:aws_eks}
\end{figure}

\section*{Conclusion}
\addcontentsline{toc}{section}{Conclusion}
Au terme de cette release, l'infrastructure de déploiement est opérationnelle : les microservices sont containerisés, validés localement sur Minikube, et l'environnement cloud AWS (EKS, RDS, ECR) est provisionné. Le chapitre suivant abordera l'automatisation du pipeline CI/CD.

\newpage
